\documentclass[11pt]{article}

\input{../shared/preamble.tex}

\usepackage{amsmath}
\usepackage{booktabs}
\usepackage{cleveref}

\title{When Does Portfolio Construction Work?\\
       A Map across the Number of Assets and the Cost of Trading}
\author{Peter Cotton}
\date{\today}

\newcommand{\Sig}{\Sigma}
\newcommand{\R}{\mathbb{R}}
\DeclareMathOperator{\diag}{diag}
\DeclareMathOperator{\tr}{tr}

\begin{document}
\maketitle

\begin{abstract}
A portfolio rule is a covariance estimate feeding an allocator. We ask a plain
question: which rules actually work out of sample---after trading costs, and as
the number of assets grows? We backtest eleven online constructions over
thousands of random subsets of U.S.\ stocks and draw a map of the winner against
two axes, the number of names and the cost of trading. The map has three regions.
With few names and cheap trading, classic minimum-variance wins. With
many names and cheap trading, a factor-model minimum-variance wins---and
by a margin that grows with the number of names---because the ordinary version
becomes undefined once there are more assets than data. Once trading
costs rise above about five to ten basis points, simple inverse-variance
weighting wins everywhere, because turnover, not estimation error, is what hurts.
We then give the two ingredients that make the many-names corner work: a Woodbury
inversion of a factor covariance, and a ridge-regularized version of the Schur
coupling.
\end{abstract}

\section{The question}
Most studies compare portfolio methods one against another and crown a winner. A
practitioner faces a different question: which method should I run, given
how many names I hold, how noisy my covariance estimate is, and how much it costs
me to trade? The answer is not a single method; it is a map.

Write a portfolio rule as
\begin{equation}
  w \;=\; \mathcal{A}\!\left(\widehat{\Sig}\right):
\end{equation}
an allocator $\mathcal{A}$ (minimum-variance, risk parity, a tilt, \dots)
applied to an estimated covariance $\widehat{\Sig}$. Whether a rule is
usable as the universe grows comes down to three things:
\begin{itemize}
  \item \textbf{Is it defined?} With more assets than data, $\widehat{\Sig}$ is
    rank-deficient and $\widehat{\Sig}^{-1}$ does not exist. Any rule that needs
    the full inverse simply has no answer.
  \item \textbf{Is it fast enough?} An $O(n^3)$ step is unusable at thousands of
    names.
  \item \textbf{Does it pay after costs?} In-sample variance is not the point;
    net-of-cost return is.
\end{itemize}
This paper measures all three at once and reports the result as a two-axis map.

\section{The methods, in one breath}
We use online estimators: a one-time \texttt{fit}, then a cheap
\texttt{partial\_fit} each period over an exponentially-weighted covariance. They
split into two families by how they use $\widehat{\Sig}$.

\paragraph{Rules that never invert the covariance (work at any size).}
Equal weight; inverse-variance ($w_i \propto 1/\widehat\sigma_i^2$,
using only the diagonal); risk parity (equal risk contributions, an
interior solution we solve by warm-started iteration)
\parencite{maillard2010}; hierarchical risk parity (HRP), which allocates
down a tree \parencite{lopezdeprado2016}; and the Thurstone tilt, which
nudges a chosen benchmark and never touches $\Sig^{-1}$.

\paragraph{Rules that invert the covariance (need it full-rank).}
Minimum-variance ($w\propto\Sig^{-1}\mathbf 1$); maximum
diversification \parencite{choueifaty2008}; maximum decorrelation; and
mean--variance. Each is a clean function of $\Sig$, but undefined when
$\Sig$ is singular---which is the normal state of affairs at scale.

\paragraph{The bridge.}
The Schur construction \parencite{cotton2024schur} has a dial $\gamma$ that slides
from HRP ($\gamma{=}0$) to minimum-variance ($\gamma{\to}1$).

\paragraph{Why turnover is built in, not bolted on.}
Because $w=\mathcal{A}(\widehat\Sig)$, weights jolt around whenever either
factor does. For the inverting rules the allocator is already a smooth function of
$\Sig$, so smooth weights just need a smooth covariance. The harder cases are HRP
and the tilt, where the allocator itself is jumpy---a re-sorted tree, a fresh
random draw---and we smooth those directly (a stable spectral ordering for the
tree; a fixed random seed reused across updates for the tilt).

\section{How we test}
We use daily returns for the U.S.\ large-cap cross-section
\parencite{winningport}; we redistribute no data. One trial works like
this:
\begin{enumerate}
  \item draw a random set of $k$ names and a random time window;
  \item walk forward---form weights from the past, earn the next day, update;
  \item record gross Sharpe, turnover, and Sharpe after a $c$-bp trading cost.
\end{enumerate}
We run thousands of trials with $k$ and the window drawn at random. Averaging over
trials removes the luck of which names and period a single backtest
happened to pick---luck large enough, we found, to flip the ranking. One more
convenience: net Sharpe is just gross minus a cost-proportional drag, so a single
run at $10$\,bp recovers the ranking at any cost.

\section{What the map shows}

\subsection{At scale, trading cost decides---not skill}
On broad subsets (hundreds of names) the gross and net winners are different
methods. Minimum-variance has the best gross Sharpe, and its raw skill even
grows with the number of names. But it trades a lot, and after costs it
falls behind. Inverse-variance and risk parity trade roughly ten times less and
win net. The crossover is early: about $5$--$10$\,bp.

\subsection{The three regions}
\Cref{tab:regime} is the headline: the winning method, by mean net Sharpe, over
$2{,}638$ trials.
\begin{table}[t]
  \centering
  \small
  \begin{tabular}{r llll}
    \toprule
    names $k\;\backslash\;$ cost & $0$\,bp & $5$\,bp & $10$\,bp & $20$\,bp \\
    \midrule
    $25$  & MinVar $0.59$ & MinVar $0.53$ & InvVar $0.50$ & InvVar $0.49$ \\
    $50$  & MinVar $0.71$ & MinVar $0.61$ & InvVar $0.56$ & InvVar $0.54$ \\
    $100$ & \textbf{FactorMV} $0.83$ & \textbf{FactorMV} $0.65$ & InvVar $0.54$ & InvVar $0.52$ \\
    $200$ & \textbf{FactorMV} $0.89$ & \textbf{FactorMV} $0.66$ & InvVar $0.56$ & InvVar $0.54$ \\
    $350$ & \textbf{FactorMV} $0.95$ & \textbf{FactorMV} $0.69$ & InvVar $0.56$ & InvVar $0.54$ \\
    $500$ & \textbf{FactorMV} $0.94$ & \textbf{FactorMV} $0.66$ & InvVar $0.53$ & InvVar $0.51$ \\
    \bottomrule
  \end{tabular}
  \caption{The regime map: winner and its mean out-of-sample net Sharpe, by number
  of names and trading cost. Dense minimum-variance (MinVar) is undefined
  past the point where names exceed the data, so it is absent there; the
  factor-model version (FactorMV) takes that corner and wins by a widening margin
  (its win rate reaches $79\%$ at $k{=}500$, zero cost).}
  \label{tab:regime}
\end{table}
Read it in three pieces:
\begin{itemize}
  \item \textbf{Few names, cheap trading $\Rightarrow$ minimum-variance.} The
    classic answer, while the covariance is still full-rank.
  \item \textbf{Many names, cheap trading $\Rightarrow$ factor minimum-variance.}
    Here ordinary minimum-variance is undefined. The factor version
    (\cref{sec:wellposed}) not only fills the gap, it wins outright, and
    its lead grows with the number of names.
  \item \textbf{Costs above ${\sim}10$\,bp $\Rightarrow$ inverse-variance,}
    everywhere. When trading is expensive, the rule that barely trades wins.
\end{itemize}
The novel coupling and tilt are competitive in the middle but rarely top the
mean; their value is staying robust and well-defined at scale, not setting the
net-Sharpe record.

\subsection{How to set the Schur dial $\gamma$}
Restricting to cases where the construction is well-defined, a middle setting of
$\gamma$ does beat HRP on raw skill---the best is near $\gamma\approx0.25$
($0.539$ versus HRP's $0.529$). But the edge is small and turnover climbs quickly
with $\gamma$, so the moment trading costs anything the best $\gamma$ drops to $0$
(plain HRP) by about $2$\,bp. The rule of thumb: pick $\gamma$ to maximise
$\text{skill}(\gamma)-\text{cost}\times\text{turnover}(\gamma)$, which at
realistic costs means keep it near zero.

\subsection{Do the shorts pay? No.}
The winning rules in the first two regions (minimum-variance and its factor
version) allow short positions. Are those shorts earning their place, or are they
just noise from inverting a noisy covariance? We test it directly: take the
allocator's own weight path, remove the shorts (clip to zero, renormalise), and
re-price it on the same returns. The shorts add value only if the signed book
beats this long-only version out of sample. On every subset we tried, it does not.
\begin{itemize}
  \item The short book grows with the number of names---its gross size reaches
    $180\%$ at $k{=}400$---which is the error-maximization one expects from an
    ill-conditioned inverse.
  \item Net of cost, dropping the shorts matches or beats keeping them at every
    dimension; only a sliver ($+0.03$ to $+0.05$ Sharpe) ever favours the shorts,
    and at large $k$ they cost $0.2$ or more.
\end{itemize}
So on this universe the negative weights are not needed: a long-only version of
the same idea does as well or better after costs. This is consistent with the
map, where the long-only inverse-variance rule already wins once trading is not
free.

\section{Making the many-names corner work}\label{sec:wellposed}
Two small fixes are what let the high-dimension column exist.

\paragraph{Factor minimum-variance (a stable inverse).}
Model the covariance as a few factors plus noise, $\Sig = BB^\top + \diag(\psi)$
with $B\in\R^{n\times k}$. Its inverse is then cheap and stable through the
Woodbury identity,
\begin{equation}
  \Sig^{-1} = D^{-1} - D^{-1}B\,(I_k + B^\top D^{-1}B)^{-1}B^\top D^{-1},
  \qquad D=\diag(\psi),
\end{equation}
costing $O(nk^2)$. Keep the noise term $\psi$ above a small floor and the inverse
always exists, no matter how many names. That single change turns
``minimum-variance is undefined'' into the winner of \cref{tab:regime}.

\paragraph{Ridge Schur (a stable coupling).}
The Schur dial needs to invert one block using another; at scale that other block
is singular. Add a small ridge to that solve. It restores well-posedness and is
self-correcting: when a block is poorly estimated the ridge dominates, the
coupling fades, and Schur slides safely back toward HRP (a large ridge reproduces
$\gamma{=}0$ exactly).

\paragraph{Scaling the tilt.}
The tilt's naive cost is $O(Mn^2)+O(n^3)$. Writing its correlation in the same
factor form makes a sample cost only $O(Mnk)$, so the benchmark-anchored,
no-inverse tilt runs at about $0.4$\,s per rebalance on $3{,}000$ names.

\section{What it all means}
One idea ties the map together. In high dimension you can only reliably estimate
simple structure: a diagonal, a few factors, or a tree. The methods that
survive are exactly the ones that use only that---inverse-variance (diagonal),
factor minimum-variance and the factor tilt (factors), HRP (tree). A rule that
demands the full $n\times n$ inverse is trying to pin down $O(n^2)$ numbers from
far less data; the only way to rescue it is to shrink it into one of those
simple forms, at which point it has effectively become one of the survivors.

The second lesson is about cost. Because the cost axis decides the winner over so
much of the map, an explicit turnover penalty and a constraint mechanism that does
not jolt the weights are not niceties---they are what let a high-skill allocator
keep its skill after costs. This sharpens the classic finding that naive
diversification is hard to beat net of cost \parencite{demiguel2009}: the one
place it is beaten is the cheap-trading, many-names corner, and the thing
that beats it is a factor optimizer, not a dense one.

\section{Conclusion}
``Is this method good?'' is the wrong question; ``when is it good?'' is the
useful one. Across the number of names and the cost of trading, the answer is a
map with three regions---dense minimum-variance, factor minimum-variance, and
inverse-variance---divided by where the covariance runs out of rank and by a
$5$--$10$\,bp cost line. At the frontier, the methods that work are the ones that
never need a full inverse, made well-defined by a factor model and a ridge and
made fast by a low-rank draw.

\printbibliography

\end{document}
